\clearpage
\chap{Extension proofs and their verification scope}\label{app:written}
The WR catalogue is deliberately separate from the eighty KS exposition items. It began with twelve complete written theorem groups. Version 1.6.0 added WR13--WR16 and supplied a focused formal counterpart for selected parts of the catalogue. Version 1.7.0 expanded all sixteen printed proofs. The accepted version 1.8.0 continuation then formalized WR01 and WR02 and checked the enlarged cumulative source environment. The present edition incorporates that evidence into the monograph. A verified pointwise estimate still does not supply a formal proof of a separate limiting conclusion.

\begin{center}
\begin{tabularx}{\linewidth}{@{}p{.19\linewidth}X@{}}
\toprule Coverage & Exact interpretation\\\midrule
Complete: \EditionWRComplete & WR01, WR02, WR06, WR07, WR13 and WR14. The principal universal statements have typed counterparts in the accepted cumulative scope. Explanatory examples retain their stated separate scope.\\[5pt]
Partial: \EditionWRPartial & WR09 and WR16. The fixed-parameter and pointwise solver/data inequalities are formal. Their parameter-limit consequences still use written WR08.\\[5pt]
Written only: \EditionWRWrittenOnly & WR03--WR05, WR08, WR10--WR12 and WR15. Full proofs are printed. No complete typed counterpart is assigned to these groups.\\\bottomrule
\end{tabularx}
\end{center}

The version 1.8.0 verification first compiled \EditionCumulativeBaseTargets\ targets from an empty project build, combining the inherited core and regularization extension. It then compiled the \EditionContinuationModules\ new WR01--WR02 modules and updated combined root over those verified objects. The fresh full audit covered \EditionCumulativeDeclarations\ declarations in the resulting environment. All \EditionCumulativePositiveFiles\ positive regression files, \EditionCumulativeCompilerNegativeFiles\ compiler rejection files, and \EditionCumulativePolicyNegativeFiles\ policy rejection files passed their specified checks. The final declaration inventory includes compiler-generated and auxiliary items. It is not a count of independent mathematical results.

The historical source files retain their original bytes. The final cumulative sources are supplied in \path{formal/}, and the accepted continuation records in \path{continuation/}. The compiler, kernel, and pinned imported dependency caches remain trusted. The source-type comparison is an editorial audit, not automatic verification of natural-language proofs. Appendix~\ref{app:concordance} distinguishes the clean baseline from the incremental continuation.

\subsection{Existence and budget geometry}
\textbf{WR01. Bounded approximate feasibility.} The real Hilbert operator and residual datum are fixed. A uniform correction-norm bound gives an exact preimage within the same bound. The printed proof uses a weakly convergent subsequence. The formal proof uses weak-star compactness of a dual ball and the Riesz isometric equivalence, with no separability or closed-range assumption. Its consequences include the range-membership equivalence and divergence of every approximating sequence's correction norms when the datum is unattained. The protected bridge applies the theorem to the actual kernel of the observable and constructs an exact repaired candidate within the original correction budget.

\textbf{WR02. Nested affine minimum corrections.} Nonempty closed affine constraint sets are nested in one real Hilbert space. The formal development constructs their minimum-norm points, proves the exact energy identity and norm monotonicity, and establishes the equivalence between a nonempty intersection and a uniform bound on the minimum norms. Under that bound, it proves strong convergence to the unique minimum-norm point of the intersection and the limiting energy identity. If the intersection is empty, the minimum norms tend to infinity. The affine hypothesis matters: the more general convex argument supplies an energy inequality rather than equality.

\textbf{WR03. Attained budget profile.} A finite nonnegative budget has an attained residual minimum even when the full operator range is not closed. The proof retains the weak topology argument rather than substituting finite-dimensional compactness.

\textbf{WR04. Sharp budget support.} The general support formula uses the actual restricted adjoint. The focused module proves the support inequality and supplies equality at a regularized optimizer's own budget. It does not prove the entire general budget-attainment route, so WR04 remains written-only as a group.

\textbf{WR05. Robust budget comparison.} The operators act on the same correction space, with operator and data errors bounded by $\eta$ and $\delta$. The bound $\delta+\eta M$ does not apply automatically when the protected kernel changes.

\subsection{Regularization and its limit}
\textbf{WR06. Constructed regularized minimizer.} Positive $\alpha$ makes $T^*T+\alpha I$ coercive. The formal proof constructs an actual continuous linear equivalence, defines the correction using its inverse, and proves the energy identity and uniqueness. The manuscript's self-adjoint closed-range proof and the formal coercive-form proof establish the same stated result. The bridge uses the existing \path{ProtectedHilbertRepair.restricted} operator.

\textbf{WR07. Parameter and budget comparison.} The formal statements compare two positive parameters, prove the two norm monotonicities, and show that no correction of smaller or equal norm has a lower residual than the regularized answer. The normalized residual has unit norm and attains the support expression whenever it is nonzero. The Lean formulation gives the optimizer and universal competitor inequality directly rather than introducing a second numerical profile definition.

\textbf{WR08. The regularization limit.} Projection onto the closure of the range is separated from actual range membership. Strong convergence of corrections requires the projected datum to be attained. Otherwise the correction norms diverge. These infinite-limit arguments remain written-only.

\textbf{WR09. Data sensitivity and parameter choice.} The focused proof supplies the coefficients $1/(2\sqrt\alpha)$ for correction noise and one for fitted-data noise. The scalar sharpness calculations and the parameter-choice convergence consequence retain their printed proofs. Formalizing a fixed-parameter inequality does not formalize its convergence application automatically.

\textbf{WR10--WR12. The explicit diagonal models.} The exact regularized family, omitted-tail bounds, matched-budget limit $\pi/4$, and attainable noisy datum remain written with exact finite enclosure tests. No general infinite-sum formalization or new diagonal Lean instance is claimed.

\subsection{Inexact solution and coordinate interpretation}
\textbf{WR13. A posteriori solver bounds.} The correction error, fitted error, and objective excess are bounded in terms of the actual normal residual and positive parameter. The focused theorem quantifies over every proposal in the Hilbert source. The protected bridge transfers its candidate-distance bound to the actual kernel subtype. Scalar sharpness examples are separated from the principal theorem and checked independently.

\textbf{WR14. A regularized primal and dual certificate.} An exact two-square identity proves weak duality, a dual maximizer constructed from the primal correction, global objective bracketing, and a correction-distance bound. The witness is arbitrary. Setting it to the proposal's fit residual gives the exact normal-residual gap.

\textbf{WR15. The coordinate covector metric.} With orthonormal ambient and residual coordinates, a full independent kernel basis gives $g=BG^{-1}q$ and $\norm g^2=q^{\mathsf T}G^{-1}q$. The invariance proof and its invertibility hypotheses are printed and rationally tested. They have not been compiled as a new matrix theorem.

\textbf{WR16. A justified stopping rule.} The focused theorem proves the pointwise sum of the solver and data bounds. The printed convergence consequence uses written WR08 and the additional assumptions $\alpha\to0$, $\delta/\sqrt\alpha\to0$, and $\epsilon/\alpha\to0$. The displayed diagonal schedule also uses written WR12.

\subsection{Public declarations for WR01 and WR02}\label{app:wr-source-map}
The following map identifies the principal clauses without reproducing the full Lean types. Those types, their actual hypotheses, and their audited axiom dependencies are recorded in \path{continuation/FORMAL_COVERAGE.json}. Every declaration below is in the accepted cumulative environment. This map does not assign separate typed coverage to every worked example or background proof.

For WR01, the declarations belong to the namespace
\begin{quote}\small\path{KomorebiSpace.BoundedApproximateRealization}.\end{quote}
\begin{longtable}{@{}p{.30\linewidth}p{.64\linewidth}@{}}
\toprule Printed clause & Declaration\\\midrule\endhead
Exact solution within the same budget & \path{exists_solution_of_bounded_residual_tendsto_zero}\\[4pt]
Range-membership equivalence & \path{mem_range_iff_bounded_approximation}\\[4pt]
Escape for unattained data & \path{norm_tendsto_atTop_of_not_mem_range}\\
\bottomrule\end{longtable}

These declarations concern bounded linear maps between real Hilbert spaces. Residual convergence is stated either as convergence of the residual norm to zero or, equivalently, convergence of the images to the fixed datum. The conclusion preserves the supplied norm bound. It does not assert strong convergence of the original approximating sequence.

For WR02, the declarations belong to the namespace
\begin{quote}\small\path{KomorebiSpace.NestedMinimumCorrections}.\end{quote}
\begin{longtable}{@{}p{.30\linewidth}p{.64\linewidth}@{}}
\toprule Printed clause & Declaration\\\midrule\endhead
Constructed minimum corrections & \path{affineCorrections_spec}\\[4pt]
Nondecreasing minimum norms & \path{nested_affine_norm_monotone}\\[4pt]
Finite-stage energy equality & \path{nested_affine_energy}\\[4pt]
Intersection and boundedness equivalence & \path{nested_affine_intersection_iff_bounded}\\[4pt]
Strong limit, unique minimum, and limiting energy equality & \path{nested_affine_strong_realization}\\[4pt]
Escape when the intersection is empty & \path{nested_affine_escape}\\
\bottomrule\end{longtable}

The Lean family is indexed by the natural numbers and is antitone under inclusion. Each affine set is explicitly nonempty and closed. Completeness is assumed for the common real Hilbert space. In the convergence theorem, the boundedness assumption concerns the constructed minimum-norm corrections themselves.

The protected application belongs to
\begin{quote}\small\path{KomorebiSpace.ProtectedBoundedRealization},\end{quote}
with public declaration
\begin{quote}\small\path{exists_exact_repair_of_bounded_approximation}.\end{quote}
It uses corrections in the actual subtype $\ker p$ with its inherited norm. For a raw candidate $f$, it constructs $y$ satisfying $Dy=0$, $py=pf$, and $\norm{f-y}\le M$. The defect's source and target are real Hilbert spaces. The observable's codomain need only be a real normed space. No closed-range hypothesis is introduced.

\subsection{Exact implementation and remaining formalization}
The retained version 1.7.0 regularization suite records \EditionRegularizationChecks\ assertions and its inexact-certificate suite \EditionInexactChecks. These are historical numerical records, not new runs performed for this edition. The independent oracle solves ambient constrained systems by determinant formulas rather than using the certificate program's kernel coordinates. The verifier checks a full kernel basis, retains its Gram metric, preserves the observable, and distinguishes measured fit from the actual residual. All outputs are exact rationals or squared bounds. It recomputes records and rejects altered claims without calling an exact optimizer. These finite tests are not universal proofs that the Python programs implement the Lean definitions.

The next formalization block is WR03--WR05, the general attained budget profile, sharp support formula, and robustness comparison. WR08, WR10--WR12, WR15, and the remaining limiting conclusions of WR09 and WR16 also retain their written status. Completing WR01 and WR02 does not automatically discharge those separate obligations.

